\documentclass[12pt]{article}

\usepackage[dvips,letterpaper,margin=0.75in,bottom=0.5in]{geometry}
\usepackage{cite}
\usepackage{slashed}
\usepackage{graphicx}
\usepackage{amsmath}
\usepackage{amssymb}
\usepackage{braket}
\usepackage{pythonhighlight}
\begin{document}

\newcommand{\ihbar}{\ensuremath{i \hbar}}
\newcommand{\Pss}{\ensuremath{\Psi^*}}
\newcommand{\dPsidt}{\ensuremath{ \frac{\partial \Psi}{\partial t} }}
\newcommand{\dPsidx}{\ensuremath{ \frac{\partial \Psi}{\partial x} }}
\newcommand{\ddPsidx}{\ensuremath{ \frac{\partial^2 \Psi}{\partial x^2} }}
\newcommand{\dPssdt}{\ensuremath{ \frac{\partial \Psi^*}{\partial t} }}
\newcommand{\dPssdx}{\ensuremath{ \frac{\partial \Psi^*}{\partial x} }}
\newcommand{\ddPssdx}{\ensuremath{ \frac{\partial^2 \Psi^*}{\partial x^2} }}

\newcommand{\dphidt}{\ensuremath{ \frac{d \phi}{dt} }}
\newcommand{\dpsidx}{\ensuremath{ \frac{d \psi}{dx} }}
\newcommand{\ddpsidx}{\ensuremath{ \frac{d^2 \psi}{dx^2} }}


\title{PHY 115L \\ Variational Methods}
\author{Michael Mulhearn}

\maketitle

\section{Introduction}

We will use a variational approach to find the ground state
$\Psi_0(x)$ of the Time-Independent Schr\"odinger Equation (TISE):
$$\hat{H} \; \psi_0(x) = E_0 \; \psi_0(x)$$
where the Hamiltonian operator is defined by:
\begin{equation}
\hat{H} = -\frac{\hbar^2}{2 m} \; \frac{d^2}{dx^2} \, + \, V(x)
\end{equation}
The variational approach is based on the observation that:
$$\braket{\psi|\hat{H}|\psi} \geq E_0$$
for all $\psi$.  Which you can see by expanding $\ket{\psi}$ in the complete set of energy eigenstates:
$$\ket{\psi} = \sum c_n \ket{E_n}$$
so:
$$\braket{\psi|\hat{H}|\psi} = \sum |c_n|^2 E_n = |c_0|^2 E_0 + \sum_{n>0} |c_n|^2 E_n$$
but:
$$\sum |c_n|^2 = 1 \;\; \implies \;\; |c_0|^2 = 1 - \sum_{n>0} |c_n|^2$$
so:
$$\braket{\psi|\hat{H}|\psi} = E_0 + \sum_{n>0} |c_n|^2 (E_n - E_0)$$
but as $E_n \geq E_0$ we have:
$$\braket{\psi|\hat{H}|\psi} \geq E_0$$
We'll explain the numerical algorithm later, but conceptually then, if
we search the Hilbert space for the lowest energy wave function, we
will find the ground state.

As before, we will use a system of units based on a characteristic length scale $a$ and an energy:
\begin{equation}
\epsilon \equiv \frac{\hbar^2}{2ma^2}
\end{equation}
With these units, we have:
\begin{equation}
\frac{\hat{H}}{\epsilon} = -a^2\frac{d^2}{dx^2} + \frac{V(x)}{\epsilon}
\end{equation}

\section{Numerical Method for Finding the Second Derivative}

Our numerical approach will rely on applying the Hamiltonian operator,
which includes a second derivative, to a trial wave function.  We will
calculate second derivatives using a second-order approximation:
$$f''(x) = \frac{f(x+h)-2f(x)+f(x-h)}{h^2}$$
We'll assume that we calculate our function with a constant spacing so that:
$$x_j = x_0 + h * j$$
and given $y$ values:
$$y_j = f(x_j)$$
we can calculate the second derivative as:
$$f''(x_j) = \frac{y_{j + 1} - 2 y_j + y_{j-1}}{h^2}$$

One challenge with this technique is how to handle the edges.  For
these exercises, it turns out that simplest solution is just to
arrange things so that the second derivative is always zero at the
boundaries.  This is the case if the edges are at $x=\pm\infty$
(effectively) or if we are considering an {\em odd} function from
$x=0$ to $x=+\infty$.  One can use periodic boundary conditions as
well, but I have found the method proposed here is a simpler solution
that works well for our exercises.

There are many ways to implement this calculation on a given numpy
array y.  A convenient and coceptually simple approach is to use the
{\tt np.roll} function, followed by setting the edges to zero by hand.
(For practice, I would prefer that you implement your own second
derivative calculation, instead of relying on a library.)

\section{The Variational Method of Finding the Ground State}

We showed in the introduction that:
$$\braket{\psi|\hat{H}|\psi} \geq E_0$$ for all $\psi$.
Our numerical algorithm will exploit this to find the ground state
through random variations of a trial wave function.

\begin{itemize}
\item Choose the range $[a,b]$ of $x$ wide enough so that the
  wave-function can be considered to be zero outside of this range,
  with the second derivative 0 at $x=a$ and $x=b$.
\item Divide $[a,b]$ into $n$ equal intervals ($n+1$ discrete points)
  with
$$x_j = a + h\,j$$ for $0 \leq j \leq n$ and step size $h = (b-a)/n$.  (This is a good case for {\tt np.linspace})
\item Choose a starting wave function $\psi_I(x)$ (an initial guess).
\item Construct an array of $y$ values from:
  $y_j = \psi_I(x_j)$
\item Set $y_0 = y_{n} = 0.$  
\item Normalize the $y$ values as described below.  The set of $y$
  values represent the initial wave function.
\end{itemize}
At each step, we normalize the wave function as follows:
\begin{itemize}  
\item Scale the $y$-values such that:
$$h \, \sum_j y_j^2 = 1$$
\end{itemize}
Starting with our initial wave function (which is automatically
accepted) we proceed iteratively, by apply a random variation to
construct a trial wave function.  We measure the energy of the trial
wave function, and accept the change if it has lower energy.
Specifically, we do the following:
\begin{itemize}
\item Choose a random integer n from $1$ to $n-1$ inclusive (notice this avoids the end points, which we are fixing to zero).
\item Choose a random (uniform) value $\delta y$ from $-\Delta y$ to $\Delta y$, where $\Delta y$ is a parameter indicating the maximum amount of change in each iteration.
\item Put $y_i \to y_i + \delta y$
\item Normalize the $y$ values as above
\item Calculate the expectation value for the energy
  $\braket{y|\hat{H}|y}$.
\item If the expectation value for the energy has decreased, accept the trial wave function as the new wave function.  If not, revert back to $y$ values before the variation from this iteration was applied.
\item Proceed to the next iteration.
\end{itemize}
We continue this process for a fixed number of iterations or until
some other criteria is established.

\section{Simulated Annealing}

A major challenge for numerical analysis is the possibility of being
stuck in a local minimum.  If the variations are too small, the
algorithm might never find the global minium (the ground state) but
instead find a local minimum that is higher than the ground state.  On
the other hand, if the variations are too large, the numerical
algorithm becomes highly inefficient, requiring a huge number of steps
to find a small variation as you become close to the maximum.

The performance of the variational method can be improved via a
process called simulated annealing.  Initially, we use a large value
of $\Delta y$, so the wave function changes rapidly, and hopefully
samples a region near the global minimum, even if we start very far
away from it.  As time goes on, we reduce the size of the $\Delta y$,
so that the wave function moves more efficiently toward the nearest
minimum.

\section{Homework Problems}

\noindent
{\bf Problem 1:} Implement a function that, when provided with an
array of $y$ values, and a step size $h$, returns an equal sized array
of values containing the second derivative at each point.  Use the
second-order numerical technique as described above and set the end
points to zero.  Validate your function for $f(x) = \sin(x)$ in the
range $[0,2\pi]$ by plotting $f(x)$ and $f''(x)$ as calculated with
your code.\\[5pt]

\noindent
{\bf Problem 2:} Implement a function {\tt norm(y,h)} that takes an
array $y$ and a step size $h$, and returns an array the same size as $y$ containeing the normalized $y$ values as described
in the write-up.  Check your function with code like this:
\begin{python}
# normalization check:
h = 0.1
y = np.random.uniform(100)
y=norm(y,h)
print(np.sum(y*y) * h)
\end{python}
which should print a number close to 1.\\[5pt]

\noindent
{\bf Problem 3:}  Implement a function {\tt energy(x,y,V)} that returns the expectation value of the energy:
$$\braket{\psi|\hat{H}|\psi}$$
The parameters $x$ and $y$ are arrays containing the discrete $x$-values ($x_j$) and the discrete $y$-values ($\psi(x_j))$ for the normalized wave function.  The parameter V is a function V(x) that defines the potential. You'll want to use your second derivative function from Problem 1 for the kinetic energy term in the Hamiltonian.  Test your energy function with code like like this:
\begin{python}
# free particle energy check:
OFFSET = 10
def V(x):
    return OFFSET
x  = np.linspace(-5,5,1000)
y  = cos(np.pi * x / 2)
y = norm(x,y)
print("calculated energy: ", energy(x,y,V))
print("expected energy:   ", OFFSET+pi**2/4)
\end{python}
Your calculated free particle energy should match the expected energy.\\[5pt]

\noindent
{\bf Problem 4:}  Implement the variational method of finding the ground state, as described above.  You'll want to use the functions you developed and tested in Problems 1-3.  Use it to find the ground state of a (nearly) infinite square well potential:
$$\frac{V(x)}{\epsilon} = \begin{cases}
0, & -a \leq x \leq a \\
B     & {\rm otherwise} \\
\end{cases}
$$
where $B$ is some big number, like $B=1000$.  Plot your initial wave function, your final wave function, and at least one intermediate wave function in the same plot. Estimate your numerical uncertainty.  Compare your results to the expected ground state wave function (inside the well):
$$\psi_0(x) = \sqrt{\frac{1}{a}}\cos\left(\frac{\pi x}{2 a}\right)$$
with ground state energy:
$$E_0 = \frac{\pi^2 \hbar^2}{2m(2a)^2} = \epsilon \, \frac{\pi^2}{4a} $$
Remember that we set $\epsilon=1$ and $a=1$ in the code.  
There are lots of possible settings, but here are the ones that I used 
(you can certainly play around with these, if you like, and see how different choices effect convergence.)
\begin{itemize}
\item I took:
$$\psi_I(x) = \begin{cases}
1, & -0.8 \leq \, x/a \, \leq 0.8 \\
0     & {\rm otherwise} \\
\end{cases}
$$
as my starting wave function.
\item I used 50 bins with $x/a$ in the range $[-1.2, 1.2]$.
\item I used simulated annealing:  (A) 10,000 iterations with $\Delta y = 0.1$, followed by (B) 10,000 iterations with $\Delta y = 0.01$, then (C) an additional 10,000 iterations with $\Delta y = 0.01$.  With these settings I reached 1\% accuracy quite rapidly.  I estimated this accuracy by the difference in energy after (B) and (C).
\end{itemize}

\noindent
{\bf Problem 5:}  Use the variational method to find the ground state of a (nearly) infinite square well potential with a curved bottom:
$$\frac{V(x)}{\epsilon} = \begin{cases}
V_p(x) & -a \leq x \leq a \\
B     & {\rm otherwise} \\
\end{cases}
$$
where
$$V_p(x) = -\epsilon \cos\left(\frac{\pi x}{2 a}\right)$$
Calculate the change in energy $\Delta E$ relative to 
$$\frac{E_0}{\epsilon} = \frac{\pi^2}{4a} $$
We can estimate this change to first order using non-degenerate perturbation theory, with
$$\Delta E = \braket{\psi_0| V_p(x) | \psi_0}$$
This amounts to the integral:
$$\Delta E = \frac{\epsilon}{a} \int_{-a}^{a} \cos^3\left(\frac{\pi x}{2 a}\right) = \epsilon \, \frac{8}{3 \pi}$$
Compare your measured energy level shift to the value expected from perturbation theory.\\

\noindent
{\bf Problem 6:}  Use the variational method to find the ground state of a harmonic oscillator:
$$V(x) = \epsilon x^2$$
Plot your initial wave function, your final wave function, and at least one intermediate wave function in the same plot. Estimate your numerical uncertainty.  Compare your ground state energy to the expected result:
$$E_0 = \epsilon$$
These are the settings I used:
\begin{itemize}
\item I used the same initial wave function as in Problem 2.
\item I used 50 bins with $x/a$ in the range $[-3, 3]$.
\item I used the same simulated annealing as in Problem 2 and reached better than $1\%$ accuracy.
\end{itemize}

\end{document}
